% =============================================================
%  Section + Appendix on the HDP-CDP error model
%  Drop-in for the main article. Numbering uses placeholders
%  (\section, etc.) and \cref-style cross-references; adapt to
%  the host paper's conventions as needed.
% =============================================================

\documentclass[11pt]{article}

\usepackage[margin=1in]{geometry}
\usepackage{amsmath, amssymb, amsthm}
\usepackage{bm}
\usepackage{mathtools}
\usepackage{hyperref}
\usepackage{natbib}

% Convenience commands (assume amsmath, amssymb, bm already loaded).
\newcommand{\E}{\mathbb{E}}
\newcommand{\Var}{\mathrm{Var}}
\newcommand{\R}{\mathbb{R}}
\newcommand{\given}{\,|\,}
\newcommand{\indep}{\perp\!\!\!\perp}

\begin{document}

% =============================================================
%  MAIN TEXT SECTION
% =============================================================

\section{Source-specific errors via a hierarchical centered Dirichlet process}
\label{sec:hdp}

The decomposition in~\eqref{eq:outcome} models the conditional mean of the log-survival time as the sum of a prognostic baseline, a treatment effect, and a source-dependent confounding bias, with the remaining residual $W_i$ absorbed into a source-specific error term $F_{S_i}$. The standard practice in the data-fusion literature is either to assume $F_0 = F_1$ outright or, more commonly, to leave the error distribution implicit by adopting a single Gaussian likelihood. Both choices are unsatisfactory in our setting. Randomized trials and real-world data sources typically differ in measurement protocol, follow-up quality, and the prevalence of extreme outcomes, and these differences plausibly enter the residual distribution rather than the structural mean. At the same time, the two error distributions are unlikely to be entirely unrelated, and a model that treats them as fully independent forfeits useful pooling, particularly when one source is small.

We address this tension with a hierarchical Dirichlet process \citep{teh2006hierarchical} prior on the source-specific mixing measures, combined with a centering construction \citep{yang2010semiparametric} that enforces the mean-zero constraint required for identifiability of the regression components. The resulting model allows $F_0$ and $F_1$ to differ flexibly while sharing information about which residual ``types'' are present in either source.

\subsection{Model}
\label{sec:hdp-model}

Conditional on the source indicator $S_i = s$, we model the residual as a location mixture of Gaussians with common scale $\sigma > 0$ and source-specific mixing measure $G_s$:
\begin{equation}
W_i \given S_i = s, G_s, \sigma \;\sim\; \int \frac{1}{\sigma}\, \phi\!\left(\frac{w - \theta}{\sigma}\right) dG_s(\theta).
\label{eq:hdp-fs}
\end{equation}
The mixing measures share latent structure through a top-level random measure $G^*$:
\begin{equation}
G^* \;\sim\; \mathrm{DP}(\gamma, H), \qquad
G_s^* \given G^*, M_s \;\sim\; \mathrm{DP}(M_s, G^*), \quad s \in \{0, 1\},
\label{eq:hdp-prior}
\end{equation}
with base distribution $H = N(0, \sigma_\theta^2)$. By Sethuraman's stick-breaking representation \citep{sethuraman1994constructive}, the top-level measure is discrete,
\begin{equation}
G^* \;=\; \sum_{k=1}^\infty \beta_k\, \delta_{\theta_k^*}, \qquad \theta_k^* \stackrel{\mathrm{iid}}{\sim} H, \qquad \beta_k = u_k \prod_{l < k}(1 - u_l), \qquad u_k \stackrel{\mathrm{iid}}{\sim} \mathrm{Beta}(1, \gamma),
\end{equation}
so that each source-specific measure $G_s^*$ is supported on the \emph{same} atoms $\{\theta_k^*\}$ as $G^*$ but carries its own weights $\bm{\pi}_s = (\pi_{sk})_{k \ge 1}$. The shape of each $F_s$ is therefore parameterized by a shared set of residual types and a source-specific weighting over them.

\subsection{Identifiability via mean-zero centering}
\label{sec:hdp-centering}

For $m_0$, $\tau$, and $c$ to retain their interpretation as the conditional-mean components defined in~\eqref{eq:cmean}, each $F_s$ must have mean zero. We follow \citet{yang2010semiparametric} and obtain centered mixing measures by subtracting the source-specific mean of the unconstrained measure:
\begin{equation}
\mu_s \;=\; \int \theta \, dG_s^*(\theta) \;=\; \sum_{k=1}^\infty \pi_{sk}\, \theta_k^*, \qquad
G_s \;=\; \sum_{k=1}^\infty \pi_{sk}\, \delta_{\theta_k^* - \mu_s}.
\label{eq:hdp-centering}
\end{equation}
By construction $\E[W_i \given S_i = s, G_s, \sigma] = 0$. The unconstrained atoms $\theta_k^*$ remain shared across sources; the centered atoms $\theta_k^* - \mu_s$ are source-specific by virtue of the source-specific shifts $\mu_0$ and $\mu_1$. The model thereby separates two distinct features of the source-specific error distributions: the \emph{locations} at which residual mass can sit, which are shared, from the \emph{prevalence} of each location and the resulting mean shift, which are source-specific.

The construction interpolates between two natural extremes. Setting $\gamma \to 0$ collapses the top-level $G^*$ onto a single atom and recovers a parametric (Gaussian) error after centering; allowing $M_0, M_1 \to 0$ with separate prior bases recovers two unrelated DP mixtures and forfeits the shared-atom structure. The HDP prior automates the trade-off through its posterior over the concentration parameters.

\subsection{Posterior computation}
\label{sec:hdp-computation}

Inference proceeds via a blocked Gibbs sampler built around a finite truncation of the stick-breaking representation at level $K$ \citep{ishwaran2001gibbs}. One sweep alternates between BART backfitting for $m_0, \tau, c$ \citep{chipman2010bart}, augmentation of censored log-survival times via truncated normal draws, conjugate updates for the unconstrained atoms $\theta_k^*$ with deterministic recentering, auxiliary-table-count augmentation for the top-level sticks \citep{teh2006hierarchical, antoniak1974mixtures}, Dirichlet-conjugate updates for the source-specific weights, and Escobar--West updates for the concentration parameters \citep{escobar1995bayesian}. Full details, including the precise form of each working residual and conditional distribution, are deferred to Appendix~\ref{app:hdp-sampler}.

The sampler preserves conjugacy with the BART updates by introducing the cluster shift $\theta_{Z_i}^* - \mu_{S_i}$ into the working residual for each forest; conditional on the latent cluster assignments, the residual model is Gaussian and the standard BART machinery applies without modification. Centering enters only as a deterministic transformation after the atom update and source-weight update, so the sampler operates entirely on the unconstrained parameterization.

\subsection{Diagnostics}
\label{sec:hdp-diagnostics}

The model provides direct posterior summaries of how much source-specific structure the error distribution carries. The posterior of $|\mu_0 - \mu_1|$ quantifies the mean discrepancy between $F_0$ and $F_1$; tight concentration near zero is evidence that a single error distribution would have sufficed, while clear separation from zero indicates that the source-specific structure is doing work. Source-specific component occupancies $n_{sk} = \sum_i \mathbf{1}\{S_i = s, Z_i = k\}$ and the auxiliary table counts $t_{s\cdot} = \sum_k t_{sk}$ produced by the sampler function as effective-sample-size-style summaries of cross-source borrowing, in the spirit of the dynamic-borrowing diagnostics of \citet{hupf2021bayesian}. The full per-source residual densities $\hat f_s(w) = \sum_k \pi_{sk}\, \phi\big((w - \theta_k^* + \mu_s)/\sigma\big)/\sigma$ can be plotted for direct visual comparison.


% =============================================================
%  APPENDIX
% =============================================================

\appendix

\section{Posterior computation for the HDP-CDP error model}
\label{app:hdp-sampler}

This appendix provides the full hierarchical model after truncation, the conditional distributions of each block of the Gibbs sampler, and recommendations for hyperparameter calibration and practical implementation.

\subsection{Truncation and full hierarchical model}
\label{app:hdp-truncated}

For computation we truncate the top-level stick-breaking at a finite level $K$ \citep{ishwaran2001gibbs}: $u_K = 1$, so $\beta_K = 1 - \sum_{k < K} \beta_k$ and $\beta_k = 0$ for $k > K$. Conditional on the truncated $\bm{\beta} \in \Delta^{K-1}$ we use the finite-dimensional approximation
\begin{equation}
\bm{\pi}_s \given \bm{\beta}, M_s \;\sim\; \mathrm{Dirichlet}(M_s \beta_1, \ldots, M_s \beta_K),
\label{eq:hdp-pi-prior}
\end{equation}
which is the standard truncated representation of the HDP \citep{teh2006hierarchical, wang2009variational}; the approximation becomes exact as $K \to \infty$. The default $K = 50$ follows \citet{henderson2020individualized}; in our simulations under $n_{\mathrm{RCT}} = 100$, $n_{\mathrm{RWD}} = 200$, the number of occupied components rarely exceeds twenty, and the largest occupied component index is monitored across iterations to confirm that the truncation does not bind.

Introducing latent cluster assignments $Z_i \in \{1, \ldots, K\}$ such that the latent location for observation $i$ is $\theta_{Z_i}^* - \mu_{S_i}$, the full hierarchical model is:
\begin{align}
Y_i \given Z_i, \bm{\theta}^*, \sigma^2, m_0, \tau, c &\sim N\!\big(\eta_i,\; \sigma^2\big), \quad \eta_i = m_0(X_i, S_i) + A_i\, \tau(X_i) + A_i(1 - S_i)\, c(X_i) + \theta_{Z_i}^* - \mu_{S_i}, \label{eq:app-lik} \\
Z_i \given \bm{\pi}_{S_i} &\sim \mathrm{Categorical}(\bm{\pi}_{S_i}), \\
\bm{\pi}_s \given \bm{\beta}, M_s &\sim \mathrm{Dirichlet}(M_s \beta_1, \ldots, M_s \beta_K), \qquad s \in \{0, 1\}, \\
\bm{\beta} \given \gamma &\sim \mathrm{Stick}_K(\gamma), \\
\theta_k^* &\stackrel{\mathrm{iid}}{\sim} N(0, \sigma_\theta^2), \qquad k = 1, \ldots, K, \\
m_0, \tau, c &\sim \mathrm{BART}, \\
\sigma^2 &\sim \kappa\nu / \chi^2_\nu, \\
\gamma &\sim \mathrm{Gamma}(a_\gamma, b_\gamma), \\
M_s &\sim \mathrm{Gamma}(a_M, b_M), \qquad s \in \{0, 1\},
\end{align}
with $\mu_s = \sum_k \pi_{sk} \theta_k^*$ a deterministic function of $\bm{\pi}_s$ and $\bm{\theta}^*$, and $\mathrm{Stick}_K(\gamma)$ the truncated stick-breaking distribution with $\mathrm{Beta}(1, \gamma)$ breaks.

\subsection{Gibbs sampler}
\label{app:hdp-gibbs}

For each individual $i$, define the partial residual
\begin{equation}
R_i \;:=\; Y_i - m_0(X_i, S_i) - A_i\, \tau(X_i) - A_i(1 - S_i)\, c(X_i).
\label{eq:app-Ri}
\end{equation}
Conditional on the BART components and the latent cluster assignment $Z_i = k$, $R_i \sim N(\theta_k^* - \mu_{S_i}, \sigma^2)$. One sweep of the sampler executes the following steps.

\paragraph{Step 1 (data augmentation for censoring).} For each $i$ with $\delta_i = 0$,
\begin{equation}
Y_i \given \cdot \;\sim\; \mathrm{TruncatedNormal}(\eta_i,\, \sigma^2;\, [\log C_i, \infty)),
\end{equation}
with $\eta_i$ as in~\eqref{eq:app-lik}.

\paragraph{Step 2 (BART backfitting).} The three forests are updated sequentially via Bayesian backfitting \citep{chipman2010bart}. The working response for each forest is the partial residual obtained by subtracting all other components and the current cluster shift $\theta_{Z_i}^* - \mu_{S_i}$. Explicitly:
\begin{align}
R_i^{(m_0)} &= Y_i - A_i \tau(X_i) - A_i(1-S_i) c(X_i) - (\theta_{Z_i}^* - \mu_{S_i}), \\
R_i^{(\tau)} &= Y_i - m_0(X_i, S_i) - (1-S_i) c(X_i) - (\theta_{Z_i}^* - \mu_{S_i}), && i \text{ with } A_i = 1, \\
R_i^{(c)} &= Y_i - m_0(X_i, 0) - \tau(X_i) - (\theta_{Z_i}^* - \mu_0), && i \text{ with } A_i = 1, S_i = 0.
\end{align}
Each working response is treated as Gaussian with mean equal to the relevant forest and variance $\sigma^2$. Control observations carry no information about $\tau$; RCT observations and untreated RWD observations carry no information about $c$. After the BART updates, recompute $R_i$ from~\eqref{eq:app-Ri}.

\paragraph{Step 3 (cluster assignments).} For each $i$,
\begin{equation}
P(Z_i = k \given \cdot) \;\propto\; \pi_{S_i, k}\, \phi\!\left(\frac{R_i - (\theta_k^* - \mu_{S_i})}{\sigma}\right), \qquad k = 1, \ldots, K.
\label{eq:app-cluster}
\end{equation}

\paragraph{Step 4 (unconstrained atoms).} Atoms are sampled in the unconstrained parameterization and the centering is reapplied deterministically. Let $n_k = \sum_i \mathbf{1}\{Z_i = k\}$ and $\bar R_k = \sum_{i: Z_i = k} R_i$. The conjugate Gaussian update is
\begin{equation}
\theta_k^* \given \cdot \;\sim\; N\!\left(\frac{\sigma^{-2}\, \bar R_k}{\sigma_\theta^{-2} + n_k\, \sigma^{-2}},\; \frac{1}{\sigma_\theta^{-2} + n_k\, \sigma^{-2}}\right), \qquad k = 1, \ldots, K.
\label{eq:app-atom-update}
\end{equation}
After updating all atoms, recompute $\mu_s = \sum_k \pi_{sk} \theta_k^*$ and $\theta_{sk} = \theta_k^* - \mu_s$ for $s \in \{0, 1\}$. The update~\eqref{eq:app-atom-update} is the standard step for the unconstrained mixing measure; centering acts as a deterministic transformation between an unidentified parameterization $(m_0^{\mathrm{unc}}, \bm{\theta}^*)$ and the identified $(m_0, \bm{\theta})$. As discussed in \citet{yang2010semiparametric} and \citet{henderson2020individualized}, this parameter-expansion strategy preserves the correct marginal posterior for the centered quantities.

\paragraph{Step 5 (top-level sticks).} Update the top-level stick weights using the auxiliary table-count augmentation of \citet{teh2006hierarchical}. Let $n_{sk} = \sum_i \mathbf{1}\{S_i = s, Z_i = k\}$. For each $(s, k)$ with $n_{sk} > 0$, sample the auxiliary table count $t_{sk}$ via the Antoniak representation \citep{antoniak1974mixtures}:
\begin{equation}
t_{sk} \;=\; \sum_{j=1}^{n_{sk}} B_{sk,j}, \qquad B_{sk,j} \;\sim\; \mathrm{Bernoulli}\!\left(\frac{M_s \beta_k}{M_s \beta_k + j - 1}\right),
\end{equation}
and set $t_{sk} = 0$ if $n_{sk} = 0$. Let $t_{\cdot k} = t_{0k} + t_{1k}$. Update the truncated sticks
\begin{equation}
u_k \given \cdot \;\sim\; \mathrm{Beta}\!\left(1 + t_{\cdot k},\; \gamma + \sum_{l > k} t_{\cdot l}\right), \qquad k = 1, \ldots, K - 1,
\end{equation}
with $u_K = 1$, and recompute $\beta_k = u_k \prod_{l < k}(1 - u_l)$.

\paragraph{Step 6 (source-specific weights).} For each $s \in \{0, 1\}$,
\begin{equation}
\bm{\pi}_s \given \cdot \;\sim\; \mathrm{Dirichlet}(M_s \beta_1 + n_{s1},\, \ldots,\, M_s \beta_K + n_{sK}).
\end{equation}
After this step, recompute $\mu_s$ and the centered atoms.

\paragraph{Step 7 (concentration parameters).} The concentration parameters $\gamma$ and $M_0, M_1$ are updated using the auxiliary-variable approach of \citet{escobar1995bayesian}. For $\gamma$, let $K^*$ denote the number of occupied top-level components and $t_{\cdot\cdot} = \sum_{s,k} t_{sk}$ the total table count. Sample $\eta_\gamma \given \gamma, t_{\cdot\cdot} \sim \mathrm{Beta}(\gamma + 1, t_{\cdot\cdot})$, then
\begin{equation}
\gamma \given \cdot \;\sim\; \omega_\gamma\, \mathrm{Gamma}(a_\gamma + K^*,\, b_\gamma - \log \eta_\gamma) + (1 - \omega_\gamma)\, \mathrm{Gamma}(a_\gamma + K^* - 1,\, b_\gamma - \log \eta_\gamma),
\end{equation}
with $\omega_\gamma / (1 - \omega_\gamma) = (a_\gamma + K^* - 1) / [t_{\cdot\cdot}(b_\gamma - \log \eta_\gamma)]$. The analogous update for $M_s$ replaces $K^*$ by the source-specific table count $t_{s\cdot} = \sum_k t_{sk}$ and uses $\eta_{M, s} \sim \mathrm{Beta}(M_s + 1, n_s)$. See \citet[Appendix A]{teh2006hierarchical} for the derivation.

\paragraph{Step 8 (error scale).} The conjugate update for the residual variance is
\begin{equation}
\sigma^2 \given \cdot \;\sim\; \mathrm{InverseGamma}\!\left(\frac{\nu + n}{2},\; \frac{\kappa\nu + \hat s^2}{2}\right), \qquad \hat s^2 = \sum_{i=1}^n \big(R_i - (\theta_{Z_i}^* - \mu_{S_i})\big)^2.
\end{equation}

\subsection{Hyperparameter calibration}
\label{app:hdp-hyperparameters}

We adopt the BART defaults of \citet{chipman2010bart} with the AFT-specific adaptations of \citet{henderson2020individualized}: $J = 200$ trees, base $\alpha = 0.95$, depth penalty $\beta = 2$, and the response transformation $y_i^{\mathrm{tr}} = y_i \exp(-\hat\mu_{\mathrm{AFT}})$ where $\hat\mu_{\mathrm{AFT}}$ is the intercept of a preliminary parametric log-normal AFT fit. The error-scale hyperparameters $\kappa$ and $\sigma_\theta^2$ are calibrated by inverting the Yamato-type approximation \citep{yamato1984characteristic} described in \citet[Section 2.4]{henderson2020individualized}: with $\sum_k \pi_{sk}(\theta_k^* - \mu_s)^2 / \sigma_\theta^2$ approximately $\chi^2_1$ in distribution, the prior on the marginal residual variance $\Var(W \given S = s)$ can be matched in tail probability to a preliminary parametric estimate $\hat \sigma_W^2$. The approximation applies within each source separately and requires no modification for the HDP construction. We use a default tail probability $q = 0.5$.

The concentration hyperpriors are $\gamma \sim \mathrm{Gamma}(a_\gamma, b_\gamma)$ and $M_s \sim \mathrm{Gamma}(a_M, b_M)$ for $s \in \{0, 1\}$, with $a_\gamma = a_M = 2$ and $b_\gamma = b_M = 0.1$. This places prior mean $20$ on each concentration with prior mode $10$, mirroring the defaults in \citet{henderson2020individualized} and the recommendations of \citet{teh2006hierarchical}.

\subsection{Practical considerations}
\label{app:hdp-practical}

\paragraph{Truncation monitoring.} Track $\max_i Z_i$ across iterations. If the largest occupied index equals $K$ in more than a negligible fraction of iterations, the truncation is too tight and $K$ should be increased. The default $K = 50$ has been adequate in all our simulations.

\paragraph{Initialization.} Initialize the BART components via the standard \citet{chipman2010bart} initialization. Assign cluster labels by $k$-means on residuals from a preliminary parametric fit; set $\bm{\beta} = (1/K, \ldots, 1/K)$, $\bm{\pi}_s = \bm{\beta}$, and $\gamma = M_0 = M_1 = 1$.

\paragraph{Posterior summaries.} Posterior draws of the centered atoms, source-specific weights, mean shifts $\mu_s$, and shared atoms $\theta_k^*$ are stored at each iteration. From these we compute (i)~per-source residual densities $\hat f_s(w)$ as in Section~\ref{sec:hdp-diagnostics}; (ii)~the posterior of $|\mu_0 - \mu_1|$; (iii)~per-source occupancy distributions $\{n_{sk}\}_k$; (iv)~the posterior of the table counts $t_{s\cdot}$, which together with $M_s$ functions as a Bayesian effective-sample-size summary of the degree of cross-source borrowing.

\paragraph{Survival-scale estimands.} The posterior of the source-$s$ survival contrasts at a covariate value $x$ is obtained by convolving the posterior draws of $m_0(\cdot, s)$, $\tau$, and $c$ with the residual mixture. For each draw, the survival function under treatment $a$ in source $s$ is $1 - F_{a, s}(\log t)$ with
\begin{equation}
F_{a, s}(y) \;=\; \sum_{k=1}^K \pi_{sk}\, \Phi\!\left(\frac{y - m_0(x, s) - a \tau(x) - a(1 - s) c(x) - \theta_k^* + \mu_s}{\sigma}\right),
\end{equation}
and the risk-difference and RMST estimands in Section~XX of the main text follow by direct evaluation. The integration step in the RMST is performed by trapezoidal quadrature over a fine grid on $[0, t^*]$.

\end{document}